\rhead{DS601: Time Series Analysis \& Forecasting}
\phantomsection 
\section*{Time Series Analysis \& Forecasting (DS601)}
\label{major:DS_6}

\noindent \small{\hyperref[main:scheme]{$\leftarrow$ Back to Scheme}} \vspace{0.5cm}

\noindent \textbf{Credits:} 3 (2-1-0)

\subsection*{Theory}
\noindent\textbf{Unit 1} \\
\textit{Time Series Fundamentals and Decomposition:} Time series components (trend, seasonality, cycle, irregular), Stationarity concepts (weak vs. strict stationarity), Trend estimation (moving averages, polynomial fitting, LOESS), Seasonal decomposition (classical, STL - Seasonal-Trend decomposition using Loess), Autocorrelation (ACF) and partial autocorrelation (PACF) analysis, Differencing and transformations for stationarity.

\vspace{0.3cm}

\noindent\textbf{Unit 2} \\
\textit{Classical Time Series Models:} ARIMA models (p,d,q parameters, stationarity/invertibility conditions), Box-Jenkins methodology, Model identification (ACF/PACF interpretation), Parameter estimation (MLE, least squares), Model diagnostics (Ljung-Box test, residual analysis), SARIMA for seasonal data, Seasonal differencing and Fourier terms.

\vspace{0.3cm}

\noindent\textbf{Unit 3} \\
\textit{Exponential Smoothing and State Space Models:} Simple exponential smoothing, Holt's linear trend method, Holt-Winters seasonal method, ETS framework (error, trend, seasonal components), Optimal smoothing parameters (MLE optimization), State space representation, Kalman filter for parameter estimation, Dynamic linear models and interventions.

\vspace{0.3cm}

\noindent\textbf{Unit 4} \\
\textit{Modern Machine Learning for Time Series:} Feature engineering (lagged variables, rolling statistics, Fourier features), Tree-based methods (XGBoost/LightGBM with time series splits), LSTM/GRU networks for sequential forecasting, Transformer models (Temporal Fusion Transformer, Informer), Ensemble methods (statistical + ML hybrid models), Cross-validation strategies (purged, embargoed, walk-forward).

\vspace{0.3cm}

\noindent\textbf{Unit 5} \\
\textit{Advanced Topics and Production Systems:} Multivariate time series (VAR, VEC models), Hierarchical forecasting (bottom-up, top-down, optimal reconciliation), Anomaly detection (isolation forest, autoencoders, statistical process control), Probabilistic forecasting (quantile regression, conformal prediction), Automated ML (AutoTS, sktime, Darts), Production deployment (MLOps for time series, concept drift detection).

\newpage

\subsection*{Practical}
\noindent\textbf{Lab 1: Time Series Exploration} \\
Focuses on decomposition, stationarity testing, ACF/PACF analysis, and trend/seasonal extraction.

\vspace{0.2cm}

\noindent\textbf{Lab 2: ARIMA Model Development} \\
Focuses on auto.arima, manual parameter selection, model diagnostics, and forecasting evaluation.

\vspace{0.2cm}

\noindent\textbf{Lab 3: Exponential Smoothing Methods} \\
Focuses on ETS model selection, Holt-Winters implementation, and smoothing parameter optimization.

\vspace{0.2cm}

\noindent\textbf{Lab 4: ML for Univariate Forecasting} \\
Focuses on feature engineering, XGBoost/LightGBM with proper time series validation.

\vspace{0.2cm}

\noindent\textbf{Lab 5: LSTM/GRU Sequence Models} \\
Focuses on multivariate LSTM networks, attention mechanisms, and early stopping.

\vspace{0.2cm}

\noindent\textbf{Lab 6: Hierarchical Forecasting} \\
Focuses on bottom-up/top-down methods and reconciliation techniques.

\vspace{0.2cm}

\noindent\textbf{Lab 7: Anomaly Detection in Time Series} \\
Focuses on statistical methods, isolation forest, and reconstruction-based approaches.

\vspace{0.2cm}

\noindent\textbf{Lab 8: Probabilistic Forecasting} \\
Focuses on quantile regression and conformal prediction intervals.

\vspace{0.2cm}

\noindent\textbf{Lab 9: Automated Forecasting Pipeline} \\
Focuses on AutoTS/Darts implementation with model selection and ensemble.

\vspace{0.2cm}

\noindent\textbf{Lab 10: Production Time Series System} \\
Focuses on complete pipeline from data ingestion through deployment and monitoring.

\vspace{0.2cm}

\newpage
\rhead{Curriculum Schema}
